% 张三慧第4版 第二卷第27章《薛定谔方程》
% 单独收集主教材原编号思考题27.1--27.7，V2 PDF物理页386（印刷369）。
% 与同号的“习题27.1--27.11”不同，绝不能重复计入705道章末习题。
\section*{第27章\quad 薛定谔方程：思考题}
\paragraph{思考题 27.1}
薛定谔方程是否由经典理论的某套严格逻辑推演唯一导出？说明其在量子力学中的基本假设地位。
\paragraph{思考题 27.2}
在含时薛定谔方程中，波函数的演化怎样保证归一化？推导对应的概率守恒条件。
\paragraph{思考题 27.3}
一个物理上可接受的波函数应满足哪些一般性质？区分连续性、有限性、单值性与归一化条件及其适用前提。
\paragraph{思考题 27.4}
波函数的归一化具有什么统计解释？为什么乘以常数不会改变归一化前的物理态，而最终概率必须规范化？
\paragraph{思考题 27.5}
比较原图27.3、27.5及27.12中的有限/无限势阱：基态概率密度在势阱内部的零点、激发态的零点数与量子数之间一般是什么关系？说明何时可使用一维节点定理。
\paragraph{思考题 27.6}
在图27.13的一维阶梯形势阱内，粒子能级为 \(E_5\)，阶梯上方仍为允许运动区域。对照 \(-a<x<0\) 与 \(0<x<a\)，哪区局部波长较长？哪区波函数振幅较大？说明近似条件。
\paragraph{思考题 27.7}
如果处于严格定态的粒子有确定能量，能否据此断定所有这种状态的寿命确定为有限值？粒子能否自发从一种严格定态跃迁到另一种？区别孤立系统与耦合电磁场的原子。
